%%%%%%%%%%%%%%%%%%%%%%%%%%%%%%%%%%%%%%%%%%%%%%%%%%%%%%%%%%%%
\section{Introduction}
\label{sec:introduction}

Enterprises are struggling to deploy agentic AI systems in production. While these systems promise to automate complex workflows—managing financial transactions, patient records, and critical infrastructure—they introduce security challenges that existing defenses cannot address. For example, within hours of releasing the OpenAI Atlas browser, researchers demonstrated that malicious instructions embedded in webpage content could trigger the assistant to exfiltrate credentials ~\cite{atlas-cve}. OpenAI's CISO acknowledged that ``prompt injection remains an unsolved problem''~\cite{openai-ciso-atlas}. 

The challenge runs deeper than prompt injection alone. LLMs cannot distinguish instructions from data. Non-deterministic execution paths cannot be predicted or statically analyzed. Multi-turn interactions allow gradual context manipulation where each message appears benign. Existing approaches look for patterns or are probabilistic—both require enumerating attacks, leaving defenders grappling with unknown unknowns.

We eliminate unknown unknowns through a systematic approach: instead of enumerating attacks, we bound the system deterministically. We build on the realization that agentic systems can be abstracted as simple, byzantine, distributed systems -- multiple entities interacting across well defined boundaries. Now security reduces to protecting boundary crossings. Conceptually, our approach is based on satisfying two properties at each boundary crossing : understanding \textbf{intent} ensuring operations satisfy organizational policies, and enforcing \textbf{integrity} to ensure operations are authentic and unmodified. Both are necessary, and neither is sufficient by itself.
This design reduces agentic workflows to a deterministic distributed system where boundaries are guarded by policies and crossings are cryptographically protected. Breaking the system requires breaking cryptography, not crafting clever prompts. Many classes of attack such as identity spoofing, session replay, policy substitution etc are eliminated by design -- they must break cryptographic primitives to succeed, others such as unauthorized data access, privilege escalation, credential exfil etc are blocked by policy at runtime.
Architecturally, \textbf{by-policy enforcement} and \textbf{by-design elimination} deliver complementary, composable defense-in-depth.

We studied 100+ agentic applications and identified key challenges: \textbf{(1) Four simultaneous attack surfaces}—tools, data, prompts, context—where attacks compose across surfaces to achieve objectives impossible through any single surface. \textbf{(2) Heterogeneous frameworks}—Agents are built using protocols (MCP, A2A), LLM interfaces (Claude, OpenAI), orchestrators (LangChain, CrewAI, AutoGen). Heterogeneity is a permanent ecosystem feature. \textbf{(3) Compositional security gaps}—per-framework checks miss violations spanning compositions. \textbf{(4) Developer burden}—requiring security logic at every interaction across heterogeneous frameworks is impractical.

Combined, the attack surface grows faster than tools can address, highlighting a critical gap: a trust layer for AI—security infrastructure positioned between application-layer defenses (don't compose) and infrastructure primitives (lack workflow semantics). 

To address (1) we introduce \textbf{authenticated workflows}—a protocol-level primitive enforcing \textit{intent} and \textit{integrity} at every interaction. We prove the 4 attack surfaces are minimal and complete. Protocol-level positioning addresses heterogeneity (2) and eliminates compositional gaps (3). Thin wrappers across 9 popular frameworks resolve developer burden (4), establishing the foundation for an enterprise AI trust layer.

\textbf{Authenticated workflows} combine cryptography, zero-trust principles, and runtime policy enforcement to ensure each operation is authentic, authorized, and attested. We present a novel design using distributed Policy Enforcement Points (PEPs) embedded at control surfaces, verifying cryptographic proofs independently with sub-millisecond overhead, requiring no centralized infrastructure. Each surface verifies operations independently before execution, creating defense in depth where compromising one surface does not compromise others.

To specify \textit{intent} we introduce MAPL, a new AI-native policy language that lets users express agentic constraints in a dynamic, scalable manner even as agents evolve and invocation context changes. MAPL provides three key capabilities: \textit{composable grammar with inheritance} enabling policies to layer and compose through intersection semantics; \textit{radically reduced specification} scaling from $O(M users \times N resources)$ to $O(\log M + N)$ through hierarchical composition and implicit principal resolution; and \textit{dynamic state via attestations}—signed claims proving operations completed, enabling policies like "export data only after anonymization-completed attestation exists" with cryptographic enforceability. 
These advances overcome limitations of traditional policy systems (OPA, AWS Cedar) in agentic environments.


Our focus is control-flow protection across operational boundaries—tool invocations and data retrievals—complementing Rajagopalan et al. \cite{rajagopalan2025prompts}, which secure data flowing through prompts and context. Together, these deliver a formal trust layer for enterprise AI. We prove practicality through a universal security runtime backing nine frameworks (Section~\ref{sec:enforcement}) via thin adapters (200–500 LOC) requiring no protocol changes—demonstrating framework-agnostic security.

Formally, we prove that the four boundaries are complete and minimal, distributed enforcement is sound, and that policy composition preserves security properties.
Empirically, we validate these properties through comprehensive attack testing and production CVE analysis.

Current agentic defenses—guardrails, prompt filters, semantic analysis (pattern matching, ML models, heuristics)—prioritize precision but deliver low or zero recall. Adversarial prompts, encoding tricks, and novel patterns routinely bypass them via semantic blind spots, while false positives block legitimate operations.
In contrast, our cryptographic enforcement guarantees high precision and high recall: operations either bear valid signatures satisfying policies or are rejected outright. Breaking the system demands shattering cryptographic primitives—not crafting clever patterns. This yields complete violation detection, zero false positives, and deterministic trust via computational hardness.


\noindent\textbf{Contributions.}
We present authenticated workflows—protocol-level primitives enforcing cryptographic verification across four control surfaces (prompts, tools, data, context), proven complete and minimal (Lemma 6). We introduce MAPL, an AI-native policy language providing cryptographic workflow dependencies via attestations and reducing policy specification from $O(M \times N)$ to $O(\log M + N)$ through hierarchical composition (Theorems 1-3). We demonstrate universal deployment across nine heterogeneous frameworks (MCP, A2A, OpenAI, Claude, LangChain, CrewAI, AutoGen, LlamaIndex, Haystack) through thin adapters (200-500 LOC) requiring zero protocol modifications, with distributed Policy Enforcement Points embedded at every boundary providing zero-trust verification with sub-millisecond overhead. Formal proofs establish completeness and soundness (Lemmas 1-7); empirical validation demonstrates deterministic guarantees (100\% recall, 0\% false positives) across 174 test cases covering 9 of 10 OWASP Top 10 risks and complete mitigation of two production CVEs (OpenAI Atlas, GitHub MCP).
